
\subsection{Character} \label{B.1}

Each electee will be assessed for exemplary character, consisting of two parts:
\begin{enumerate}
    \item \subsecitem{a}{Electee Questionnaire} Electees will submit written, free response answers, including one long essay related to exemplary character. 
    Questions will be determined by the Vice Presidents with input from the officers.
    \item \subsecitem{b}{Interview} The electee must undergo an interview conducted by at least one active member. The Vice Presidents will determine the length and scope of the interview. 
    It is recommended that the interviewer(s) are also graduate students. 
\end{enumerate}


\subsection{Meetings} \label{B.2}

Each electee is required to attend all meetings as specified in Bylaw \ref{V.1} that are open to all members of the chapter, with the exception of the 
candidate informational meeting if they have demonstrated interest in the election process. Given the typical length of the Officer Elections Meeting, 
electees who stay longer than two hours will be awarded service hours in one-hour increments for the remainder of their attendance. 


\subsection{Service} \label{B.3}

Each electee must complete a total of 10 service hours through service projects offered by the chapter. 
Electees may report up to 5 hours of external (non-TBP sponsored) service towards this requirement at the discretion of the Service Coordinator.

\subsection{Socials and Professional Development} \label{B.4}

Each electee must attend at least three social or professional development events sponsored by Tau Beta Pi. 
At least one of these social events must be a graduate student social as designated by the Graduate Student Vice President.

\subsection{Dues} \label{B.5}

Each electee must pay the one-time-only membership dues as set by the Officer Corps. 
Any electee who is unable to pay the fee may submit a chapter fee waiver form or apply for a National Initiation-Fee Loan, and will be approved at the discretion of the Treasurer.

%Added F26
\subsection{Electee Exam} \label{B.6}
Each electee must complete an electee exam that is written by the Vice Presidents with input from the officers. 
The exam must include questions related to the National Constitution and Bylaws, in accordance with the \href{http://www.tbp.org/off/ConstBylaw.pdf}{Tau Beta Pi Association Constitution Article III Section 4(b)}. 


\subsection{Group Meetings} \label{B.7}

Each electee must participate in one group activity with their electee group. 
This can be in the form of a social or service activity but does not have to be a Tau Beta Pi sponsored event.

\subsection{Advisor Form} \label{B.8}

Each electee must get the standard advisor recommendation
form signed by their advisor or graduate program coordinator and turned into the
Graduate Student Vice President by their interview.

\subsection{Initiation} \label{B.9}

Each electee must attend initiation at the end of the term. This
is an absolute requirement and cannot be made up. An electee who has completed all other
requirements but who misses initiation may not become a member of Tau Beta Pi. Such an
electee may either attend the initiation of another chapter or must wait until the following
term to attend the initiation.


\subsection{Timing and Deadlines} \label{B.10}

The Vice Presidents, with input from the Officer Corps, will establish appropriate requirement deadlines and substitution policies each semester. 
An electee may choose to count any activities completed after initiation but before the end of the term toward the following semester, or, 
if approved by the Graduate Student Vice President, the current term.